% Generated draft: requires mathematical and source review.
\documentclass[11pt]{article}
\usepackage[margin=25mm]{geometry}
\usepackage{fontspec}
\setmainfont{Latin Modern Roman}
\usepackage{amsmath,amssymb,mathtools,unicode-math}
\setmathfont{Latin Modern Math}
\usepackage{enumitem,needspace,xurl,hyperref}
\setlength{\parindent}{0pt}
\setlength{\parskip}{5pt}
\setlength{\emergencystretch}{3em}
\newcommand{\N}{\mathbb{N}}
\newcommand{\Z}{\mathbb{Z}}
\newcommand{\Q}{\mathbb{Q}}
\newcommand{\R}{\mathbb{R}}
\newcommand{\C}{\mathbb{C}}
\newcommand{\F}{\mathbb{F}}
\newcommand{\A}{\mathbb{A}}
\newcommand{\PP}{\mathbb{P}}
\newcommand{\E}{\mathbb{E}}
\newcommand{\eps}{\varepsilon}
\newcommand{\dd}{\,\mathrm{d}}
\begin{document}
\section*{UnsolvedMath ID: 30003458}
\textbf{Source problem code:} catalogue entry
\subsection{Definitions and mathematical statement}
The supplied record asks for clarification of the displayed decomposition
\[
\left.\mathrm{Sht}^{\mu}_n\right|_{r=0}
=\coprod_E\bigl[\operatorname{Spec}\mathbb{F}_q/\operatorname{Aut}(E)\bigr].
\]
Here \(\operatorname{Aut}(E)\) denotes the automorphism group of the object \(E\), with its category and any scheme or sheaf structure left unspecified by the record. The notation \(\mathbb{F}_q\) denotes a finite field of cardinality \(q\), and thus presupposes that \(q\) is a prime power.

For terminology only, once a group sheaf \(G\) over a base \(S\) and a topology \(\tau\) have been specified, a right \(G\)-torsor on \(S\) is a sheaf \(P\) with a right \(G\)-action that is locally isomorphic, in topology \(\tau\), to \(G\) with its right translation action. Neither this topology nor the intended nature of \(\operatorname{Aut}(E)\) is supplied here.

\emph{Recorded clarification question.} For \(r=0\), how should the displayed decomposition be interpreted concretely as describing \(\operatorname{Aut}(E)\)-torsors on \(S\)?

The record does not define the shtuka stack \(\mathrm{Sht}^{\mu}_n\), the parameters \(n\) and \(\mu\), the objects indexing \(E\), the permitted bases \(S\), or the precise meaning of the displayed equality. Consequently, it does not determine a standalone mathematical proposition with complete hypotheses and quantifiers. Its supplied verification note says that the report immediately answers the clarification using an automorphism-group torsor and a twisted form. This is recorded source commentary, not an independently verified assertion here.
\subsection{Short English statement}
The source asks how the \(r=0\) shtuka decomposition describes automorphism-group torsors on a base \(S\). Essential definitions and hypotheses are absent. Although the registry labels the item open, its supplied verification note identifies it as an answered clarification question, so no precise open conjecture can be extracted.
\subsection{Sources}
\begin{itemize}
\item Local UnsolvedMath registry, permanent ID \(30003458\), exact source problem code \texttt{catalogue entry}, title ``Shtuka Strata as Automorphism Torsors.'' No individual proposer is supplied; the proposed year is recorded as 2017.
\item Supplied report citation: \emph{Arbeitsgemeinschaft: Higher Gross Zagier Formulas (2018)}, Oberwolfach Reports, pp.~1067--1134, DOI \texttt{10.4171/owr/2017/18}. Supplied source URL and literature reference ``Original Oberwolfach report'': \url{https://doi.org/10.4171/owr/2017/18}. The report itself was not consulted for this draft.
\item Supplied literature assessment, checked on exactly \texttt{2026-08-21}: the extracted item is described as the clarification question ``What exactly does this mean?'' and retained as open under a triage rule because it lacks a standalone proposition. This is supplied dated curation, not an independently verified literature assessment.
\item Supplied statement verification: the official report is said to explain the stratum immediately as an automorphism-group torsor and a twisted form, making this an answered clarification rather than an open problem. The record marks the statement \texttt{unrecoverable} and promotes no clean statement; the original remains visible for provenance.
\end{itemize}
\end{document}
